(1) 内球全部向外发射的光最终被外球内壁吸收，因为外球包围内球；内球面积为 \(A_1=4\pi r^2\)。黑体间的净辐射传热功率由两方向的辐射平衡得到
\[
Q_{1\to2}=\sigma A_1(T_1^4-T_2^4).
\]
这里外球表面积更大，但内球只截获相应几何份额的外球内壁辐射，因此净交换面积仍为内球面积。

(2) 两球除辐射外无外界热量或功交换，故
\[
C_1\dot T_1=-Q_{1\to2},\qquad
C_2\dot T_2=Q_{1\to2}.
\]
直接相加得 \(d(C_1T_1+C_2T_2)/dt=0\)。若 C₁,C₂ 温度无关，最终共同温度由能量守恒决定为 \(T_f=(C_1T_{1i}+C_2T_{2i})/(C_1+C_2)\)。

(3) 热容恒定时各球熵微分为 \(dS_i=C_i\,dT_i/T_i\)，故
\[
\dot S_{\rm total}=-\frac Q{T_1}+\frac Q{T_2}
=\sigma A_1(T_1^4-T_2^4)\left(\frac1{T_2}-\frac1{T_1}\right)\ge0.
\]
因 \(T_1^4-T_2^4\) 与 \(T_1-T_2\) 同号，而 \(1/T_2-1/T_1=(T_1-T_2)/(T_1T_2)\) 亦同号，所以总熵非负，等温时恰为零。

(4) 令 \(T_1=T_0+\delta T_1,T_2=T_0+\delta T_2\)，一阶近似 \(T_1^4-T_2^4\simeq4T_0^3(\delta T_1-\delta T_2)\)。定义温差 \(\Delta=T_1-T_2\)，联立得到
\[
\dot\Delta=-4\sigma A_1T_0^3\left(\frac1{C_1}+\frac1{C_2}\right)\Delta.
\]
故 \(\Delta(t)=\Delta(0)e^{-t/\tau}\)，时间尺度 \(\tau=[4\sigma A_1T_0^3(C_1^{-1}+C_2^{-1})]^{-1}\)，热容越大或辐射换热越弱，温差消失越慢。
